%! AP Statistics — Unit 2, Lesson 2.9 — Homework
\documentclass[10pt]{article}
\usepackage{apstats-article}
\usepackage{apstats-boxes}

\begin{document}

\pageheader{Unit 2, Lesson 2.9}{Homework}
\namedateperiod

\vspace{0.18in}

% ── DIRECTIONS ───────────────────────────────────────────────────────────────
\noindent\textbf{Directions:} Show all work. Write equations using variable names. For every
prediction, include the back-transformation step.

\vspace{0.14in}

% ── PROBLEM 1 ─────────────────────────────────────────────────────────────────
\noindent\textbf{1.}\enspace A forester measured the age (years, $x$) and diameter at breast
height (cm, $y$) of 30 Douglas fir trees. The scatterplot showed a curved upward pattern.
After applying $\log_{10}$ to $y$, the new scatterplot appeared linear. Technology output
for the regression of $\log(\text{diameter})$ on age:
\[
  \widehat{\log(\text{diameter})} = 0.62 + 0.019x \qquad r^2 = 0.95
\]

\begin{enumerate}[label=\textbf{(\alph*)}, leftmargin=2em, itemsep=0.18in]

  \item Which variable was transformed? What type of model does this suggest?\\[8pt]
        \writeline\\[3pt]\writeline

  \item Write the back-transformed equation for predicting diameter (in original units).
        Show your work.\\[8pt]
        \writeline\\[3pt]\writeline\\[3pt]\writeline

  \item Predict the diameter of a tree that is \textbf{40 years old}. Show all steps.\\[8pt]
        \writeline\\[3pt]\writeline\\[3pt]\writeline

  \item The original linear fit had $r^2 = 0.74$. Compare this to the transformed model's
        $r^2 = 0.95$. What does this comparison tell you?\\[8pt]
        \writeline\\[3pt]\writeline\\[3pt]\writeline

\end{enumerate}

\vspace{0.18in}

% ── PROBLEM 2 ─────────────────────────────────────────────────────────────────
\noindent\textbf{2.}\enspace A physicist measured the intensity ($y$, watts/m$^2$) of light
at various distances from a source ($x$, meters). Both $x$ and $y$ were $\log_{10}$
transformed. Technology output:
\[
  \widehat{\log(\text{intensity})} = 2.01 - 1.98\,\log(\text{distance}) \qquad r^2 = 0.999
\]

\begin{enumerate}[label=\textbf{(\alph*)}, leftmargin=2em, itemsep=0.18in]

  \item Which variable(s) were transformed? What type of model is suggested?\\[8pt]
        \writeline\\[3pt]\writeline

  \item Write the back-transformed power-model equation. Show work.\\[8pt]
        \writeline\\[3pt]\writeline\\[3pt]\writeline

  \item Predict the intensity at a distance of \textbf{5 meters}. Show all steps.\\[8pt]
        \writeline\\[3pt]\writeline\\[3pt]\writeline

  \item The theoretical inverse-square law gives $y \propto x^{-2}$. Does the data support
        this? Explain using the regression output.\\[8pt]
        \writeline\\[3pt]\writeline\\[3pt]\writeline

\end{enumerate}

\vspace{0.18in}

% ── PROBLEM 3 (AP-STYLE) ──────────────────────────────────────────────────────
\noindent\textbf{3.}\enspace Computer output for two models fit to the same data is shown.

\vspace{8pt}
\renewcommand{\arraystretch}{1.4}
\begin{tabularx}{\linewidth}{@{} l X X @{}}
\rowcolor{sky} \textbf{Model} & \textbf{Equation} & \textbf{$r^2$} \\
A (linear) & $\hat{y} = 3.5 + 2.8x$ & 0.82 \\
B (log $y$ on $x$) & $\widehat{\log y} = 0.72 + 0.14x$ & 0.96 \\
\end{tabularx}

\vspace{10pt}
The residual plot for Model A shows a curved pattern. The residual plot for Model B shows
random scatter.

\begin{enumerate}[label=\textbf{(\alph*)}, leftmargin=2em, itemsep=0.18in]

  \item Which model is more appropriate for these data? Justify using two pieces of evidence.\\[8pt]
        \writeline\\[3pt]\writeline\\[3pt]\writeline

  \item Using Model B, predict $y$ when $x = 5$. Show all work including back-transformation.\\[8pt]
        \writeline\\[3pt]\writeline\\[3pt]\writeline

\end{enumerate}

\vspace{0.18in}

% ── EXTENSION ─────────────────────────────────────────────────────────────────
\begin{extensionbox}
\small\textbf{Extension --- Optional}

Research one example from science, economics, or medicine where a logarithmic or power
transformation is routinely used (e.g., the Richter scale, decibels, the pH scale,
population growth models). Write: (1) what variable(s) are transformed and why, (2) the
general form of the resulting model, and (3) one real prediction the model makes. Cite
your source.
\end{extensionbox}

\vspace{0.16in}

% ── PREVIEW ───────────────────────────────────────────────────────────────────
\begin{tcolorbox}[colback=greenbg, colframe=greenacc, boxrule=0.8pt, arc=2mm,
    left=4mm, right=4mm, top=3mm, bottom=3mm]
\small\textbf{Preview --- Unit 3}\\[4pt]
Unit 2 has been about \emph{analyzing} two-variable data. Unit 3 shifts to \emph{collecting}
data. The key question: does the way we collect data affect whether we can generalize results
to a population or establish cause-and-effect? We'll explore \textbf{sampling methods} and
\textbf{experimental design}.\\[6pt]
\textbf{Think ahead:} A researcher wants to know whether a new fertilizer increases tomato
yield. Should they observe existing farms, or design an experiment? What's the difference?\\[6pt]
\writeline\\[2pt]\writeline
\end{tcolorbox}

\end{document}
